% 可直接插入中文数学建模论文正文。
% 导言区至少需要：
% \usepackage{amsmath,amssymb,bm,booktabs,threeparttable}
% \usepackage{algorithm,algpseudocode}
% \usepackage{tikz}
% \usetikzlibrary{arrows.meta,positioning,shapes.geometric}

\section{第二监测点的边界感知Minimax选址模型}
\label{sec:q2-minimax}

\subsection{第一次观测后的集合表示}

设目标圆域为
\begin{equation}
 \mathcal D=\{\bm x\in\mathbb R^2:\|\bm x\|_2\le1800\},
\end{equation}
第一监测点和示向度分别为 $\bm s_1$ 和 $\hat\theta_1$，测向半误差为
$\delta=1^\circ$。记 $\mathcal W(\bm s,\theta,\delta)$ 为以 $\bm s$ 为顶点
$\theta$ 为中心方向、$\delta$ 为半角的前向角域。第一次返回示向度后，干扰源位置可行域为
\begin{equation}
 \mathcal X_1=\mathcal D\cap B(\bm s_1,1500)
 \cap\mathcal W(\bm s_1,\hat\theta_1,\delta).
 \label{eq:q2-first-set}
\end{equation}
其中 $B(\bm c,r)$ 表示闭圆盘。目标圆与窄角域相交后可能产生圆弧边界，故
$\mathcal X_1$ 一般不是规则扇形。鉴于 5 m 近距离区相对于 1000--1500 m 接收尺度极小，
主模型暂忽略该局部返回差异，其影响在文末单独讨论。

\subsection{不遗漏保证接收点的完整候选域}

对任意 $\bm x\in\mathcal X_1$，令
\begin{equation}
 \rho_{\min}(\bm x)=\max\{1000,\|\bm x-\bm s_1\|_2\}.
\end{equation}
这是与第一次成功接收相容的最小有效接收半径。定义第二点的可靠接收域
\begin{equation}
 \boxed{\mathcal C_{\rm rec}
 =\bigcap_{\bm x\in\mathcal X_1}
 B\!\left(\bm x,\rho_{\min}(\bm x)\right).}
 \label{eq:q2-complete-safe}
\end{equation}
若 $\bm q\in\mathcal C_{\rm rec}$，则对任一相容隐藏半径 $\rho$ 都有
$\|\bm q-\bm x\|\le\rho_{\min}(\bm x)\le\rho$，第二点必能接收信号。反之，
若 $\bm q\notin\mathcal C_{\rm rec}$，则存在 $\bm x\in\mathcal X_1$ 使
$\|\bm q-\bm x\|>\rho_{\min}(\bm x)$；取
$\rho=\rho_{\min}(\bm x)$ 即得到一个第一次可接收、第二次不可接收的相容状态。
因此式~\eqref{eq:q2-complete-safe} 是保证接收的充要条件，不会遗漏任何保证接收点。

为计算式~\eqref{eq:q2-complete-safe}，令
\begin{equation}
 \bm x=\bm s_1+r\bm u(\phi),\qquad
 \phi\in[\hat\theta_1-\delta,\hat\theta_1+\delta],
\end{equation}
其中 $\bm u(\phi)=(\cos\phi,\sin\phi)^{\mathsf T}$。射线与目标圆的交点参数为
\begin{equation}
 r_{\pm}(\phi)=-\bm s_1^{\mathsf T}\bm u(\phi)
 \pm\sqrt{(\bm s_1^{\mathsf T}\bm u(\phi))^2+1800^2-\|\bm s_1\|_2^2}.
 \label{eq:q2-ray-circle}
\end{equation}
故该方向的实际径向区间为
\begin{equation}
 \ell(\phi)=\max\{0,r_-(\phi)\},\qquad
 L(\phi)=\min\{1500,r_+(\phi)\}.
\end{equation}
当 $r\le1000$ 时，$\|\bm q-\bm s_1-r\bm u(\phi)\|^2$ 关于 $r$ 为凸二次函数，
只需检查径向区间的两个端点；当 $r\ge1000$ 时，将接收条件平方可得
\begin{equation}
 \|\bm q-\bm s_1\|^2
 \le2r(\bm q-\bm s_1)^{\mathsf T}\bm u(\phi),
\end{equation}
其关于 $r$ 为线性约束，只需检查远距离区间的最小端点。于是二维连续状态检查被降为
$2^\circ$ 小区间上的一维角度搜索，并自然包含向外裁切和相切情形。

\subsection{最小覆盖圆指标及其求法}

对凸多边形 $P=\operatorname{conv}\{\bm v_1,\ldots,\bm v_n\}$，定义
\begin{equation}
 R_{\rm MEC}(P)=\min_{\bm c}\max_i\|\bm v_i-\bm c\|_2.
 \label{eq:q2-mec}
\end{equation}
任意圆盘包含全部顶点当且仅当包含其凸包，因此只需处理顶点。最小覆盖圆由两个或三个
边界顶点支撑：两点情形中二者为圆的直径端点；三点情形中圆为三点外接圆，且不可约三点
构成锐角三角形。由此可枚举全部顶点对和非共线三元组，保留覆盖所有顶点的候选圆，并取
半径最小者。

需要指出，仅先求一组全局直径端点，再固定该点对与其他顶点构圆并取最大值，一般不能得到
最小覆盖圆。正确的快速分支是：先检查直径圆；若其覆盖全部顶点，则由
$R_{\rm MEC}\ge\operatorname{diam}(P)/2$ 可知它已经最优；否则必须允许支撑点改变，并执行完整
两点、三点枚举。

圆弧区域采用内外正多边形逼近
\begin{equation}
 P_N^-\subseteq P_{\rm exact}\subseteq P_N^+,
\end{equation}
从而有
\begin{equation}
 R_{\rm MEC}(P_N^-)\le R_{\rm MEC}(P_{\rm exact})
 \le R_{\rm MEC}(P_N^+).
\end{equation}
外接正 $N$ 边形的最大径向超差为
$R(\sec(\pi/N)-1)$；当 $N=512$、$R=1800$ m 时该值为 $0.0339$ m。

\subsection{条件Minimax策略和候选区域}

给定第二点 $\bm q$ 和可能第二示向度 $z$，两次观测后的定位区域为
\begin{equation}
 \mathcal P_2(\bm q,z)=\mathcal X_1\cap B(\bm q,1500)
 \cap\mathcal W(\bm q,z,\delta).
 \label{eq:q2-second-set}
\end{equation}
定义
\begin{equation}
 J_R(\bm q)=\sup_{z\in\Theta(\bm q)}
 R_{\rm MEC}(\mathcal P_2(\bm q,z)).
 \label{eq:q2-minimax-radius}
\end{equation}
第二监测点为
\begin{equation}
 \boxed{\bm q^*\in\arg\min_{\bm q\in\mathcal C_{\rm rec}}J_R(\bm q).}
 \label{eq:q2-decision}
\end{equation}
在半径同分时，依次以最坏面积和移动距离打破同分。

原题未规定数值精度门槛，因此采用 5\% 近优集合回答“候选区域”：
\begin{equation}
 \mathcal C_{5\%}=\{\bm q\in\mathcal C_{\rm rec}:
 J_R(\bm q)\le1.05J_R(\bm q^*)\}.
 \label{eq:q2-near-optimal}
\end{equation}
实现中为接近零半径的情形保留 $0.05$ m 绝对容差，即将右端阈值取为
$\max\{1.05J_R(\bm q^*),J_R(\bm q^*)+0.05\}$。
若与后续 20 m 清除半径衔接，可另给应用型集合
\begin{equation}
 \mathcal C_{20}=\{\bm q\in\mathcal C_{\rm rec}:J_R(\bm q)\le20\}.
\end{equation}
当 $\mathcal C_{20}$ 为空时仍输出 $\mathcal C_{5\%}$，不能据此认为问题二无可用策略。

\begin{algorithm}[htbp]
\caption{边界感知Minimax MEC第二监测点算法}
\label{alg:q2-minimax-mec}
\begin{algorithmic}[1]
\Require $\bm s_1,\hat\theta_1,\delta$，圆域离散边数 $N$
\Ensure $\bm q^*$、$\mathcal C_{5\%}$ 和可选的 $\mathcal C_{20}$
\State 构造第一次可行域 $\mathcal X_1$
\State 由式~\eqref{eq:q2-ray-circle} 计算各方向真实径向区间
\State 利用径向端点约化构造完整可靠接收域 $\mathcal C_{\rm rec}$
\State 在 $\bm s_1$ 周围生成全角度粗网格并筛选可靠候选点
\For{每个 $\bm q\in\mathcal C_{\rm rec}$}
  \State 枚举几何临界方向和离散未来示向度 $z$
  \State 裁剪得到 $\mathcal P_2(\bm q,z)$
  \State 由两点、三点支撑圆计算 $R_{\rm MEC}(\mathcal P_2)$
  \State 更新最坏示向度下的 $J_R(\bm q)$
\EndFor
\State 在粗网格最优点附近局部加密并重复评分
\State 输出 $\bm q^*$、$\mathcal C_{5\%}$ 和 $\mathcal C_{20}$
\end{algorithmic}
\end{algorithm}

\subsection{仿真结果}

全局实验包含 48 组配对随机案例，另设置 16 组第一点位于半径 1500--1775 m 且测向朝外的
边界压力案例。各策略使用完全相同的真值、接收半径和误差样本。概率分布仅用于仿真比较，
不作为Minimax模型的理论前提。

\begin{table}[htbp]
\centering
\caption{全局随机案例的MEC半径比较}
\label{tab:q2-global-mec}
\small
\begin{tabular}{lrrrrr}
\toprule
策略 & 平均值/m & P95/m & 最大值/m & 平均条件最坏值/m & $P(R\le20)$\\
\midrule
本文Minimax MEC & \textbf{23.22} & \textbf{41.59} & \textbf{52.30} & \textbf{47.51} & 37.5\%\\
远场解析 & 26.13 & 42.50 & 55.65 & 54.46 & 35.4\%\\
固定$(750,375)$ & 25.79 & 59.25 & 75.36 & 66.09 & \textbf{54.2\%}\\
\bottomrule
\end{tabular}
\end{table}

配对差分表明，本文策略相对远场解析平均降低实际MEC半径 2.91 m，95\% 区间为
$[1.18,4.64]$ m；平均条件最坏半径降低 6.95 m，区间为 $[4.79,9.12]$ m。
相对固定点，上述两项分别降低 2.56 m 和 18.58 m。固定点的 20 m 达标率较高但尾部显著更差，
这是因为本文优化最坏半径而非达标概率，两种目标不可混写。

\begin{table}[htbp]
\centering
\caption{边界向外压力案例的MEC半径比较}
\label{tab:q2-edge-mec}
\small
\begin{tabular}{lrrrrr}
\toprule
策略 & 平均值/m & P95/m & 最大值/m & 平均条件最坏值/m & $P(R\le20)$\\
\midrule
本文Minimax MEC & \textbf{5.35} & \textbf{8.85} & \textbf{10.72} & \textbf{6.69} & \textbf{100\%}\\
远场解析 & 34.74 & 42.73 & 43.48 & 43.50 & 12.5\%\\
固定$(750,375)$ & 25.33 & 31.29 & 31.33 & 31.62 & 12.5\%\\
\bottomrule
\end{tabular}
\end{table}

在边界压力案例中，本文策略相对远场解析和固定点分别降低平均MEC半径 29.39 m 和
19.97 m。其原因是边界感知模型利用了第一次可行域被目标圆截短的信息，并允许在旧解析
safe 内集之外选择仍具有确定接收保证的点。

\subsection{数值可靠性与局限}

代表点将圆边数由 64 增至 512、未来示向角步长由 $2^\circ$ 缩至 $0.5^\circ$ 时，中心案例
最坏MEC半径变化 0.0191 m，边界向外案例在显示精度内不变。可靠接收域角度步长由
$0.2^\circ$ 缩至 $0.0125^\circ$ 时，测试点安全裕度保持稳定。所有 64 组案例的真值包含率和
离散可靠接收约束通过率均为 100\%。当前粗网格加局部加密算法给出的是规定分辨率下的数值Minimax解，
不是连续动作域的全局最优证明；候选区域也应连同网格分辨率一并报告。

题目中的 5 m 局部返回条件仅占最小接收半径的 $5/1000=0.5\%$，对应圆盘面积只占
目标圆面积的 $(5/1800)^2=7.72\times10^{-6}$。因此本文在第二点策略的主模型中忽略该条件；
本次 64 个案例、3 种策略共形成 192 次执行，第二点落入真值 5 m 邻域的次数为 0，
故这一简化没有改变本次实验指标。
